\documentclass[10pt, oneside]{article}   	% use "amsart" instead of "article" for AMSLaTeX format
\usepackage{geometry}                		% See geometry.pdf to learn the layout options. There are lots.
\geometry{letterpaper}                   		% ... or a4paper or a5paper or ... 
%\geometry{landscape}                		% Activate for rotated page geometry
%\usepackage[parfill]{parskip}    		% Activate to begin paragraphs with an empty line rather than an indent
\usepackage{graphicx}				% Use pdf, png, jpg, or eps§ with pdflatex; use eps in DVI mode
								% TeX will automatically convert eps --> pdf in pdflatex		
\usepackage{amssymb}
\usepackage{mhchem}
\usepackage{hyperref}

\title{Module 15, Part 2}
\author{Mark Peever\\ \texttt{mpeever@gmail.com}}
\date{April 21, 2026}

\begin{document}
\maketitle

\section*{Objectives}
\marginpar{0 minutes}
Refer to \href{https://drive.google.com/file/d/1_Hv3SaslAcLJrZ5UijrZkMslhQ-wCpJq/view?usp=sharing}{Module 15 Notes}.\\

By the end of this class, the students should be able to\ldots
\begin{itemize}
\item describe \textbf{Acid/Base Equilibrium}
\item describe the \textbf{pH Scale}
\end{itemize}

\section*{Welcome \& Devotion}
\marginpar{5 minutes}
\begin{itemize}
\item have one student read \href{https://www.biblegateway.com/passage/?search=Ecclesiastes\%209\%3A10–18\&version=LSB}{Ecclesiastes 9:10--18}
\end{itemize}

\section*{Module 14 Quiz}
\marginpar{20 minutes}



\section*{Acid/Base Equilibrium}
\marginpar{20 minutes}

\begin{itemize}
\item Acid/Base reactions are special cases of Chemical Equilibrium
\item Acid Ionization Reaction: 
\begin{enumerate}
\item \ce{HCl (aq)  <=> H+ (aq) + Cl- (aq)} 
\item \ce{HCl (aq) + H2O (l) <=> H3O+ + Cl- (aq)}
\end{enumerate}

\item Base Ionization Reaction:
\begin{enumerate}
\item \ce{NH3 (aq)  + H2O (l) <=> NH4+ (aq) + OH- (aq)} 
\item \ce{NaOH (aq)  <=> Na+ (aq) + OH- (aq)}
\end{enumerate}
 
\item introduce Acid Ionization Constant($K_{a}$) 
\item introduce Base Ionization Constant($K_{b}$) 
\end{itemize}

\section*{The pH Scale}
\marginpar{40 minutes}
\begin{itemize}
\item describe $pH$
\item $pH = - \log{10} [H^{+}]$
\item spend \emph{lots} of time on \textbf{logarithms}
\item $ 0 \le pH \le 14$
\item in other words,  $1 \times 10^{-14} M \le [ H+ ] \le 1 \times 10^{0} M$
\item here's the twist: the \emph{lower} the $pH$, the \emph{more} acidic a solution is
\item here's the twist: the \emph{higher} the $pH$, the \emph{more} basic (alkaline) a solution is
\end{itemize}

\section*{Questions for me}
\marginpar{5 minutes}

\section*{Assignment}
\marginpar{5 minutes}
\begin{itemize}
\item Review Problems: p. 522 \# 1--10 (not to be turned in)
\item Practice Problems: pp. 523--524 \# 1--10 (due 2026-05-01)
\end{itemize}



\end{document}  